%%%%%%%%%%%%%%%%%%%%%%%%%%%%%%%%%%%%%%%%%
% Developer CV
% LaTeX Template
% Version 1.0 (28/1/19)
%
% This template originates from:
% http://www.LaTeXTemplates.com
%
% Authors:
% Jan Vorisek (jan@vorisek.me)
% Based on a template by Jan Küster (info@jankuester.com)
% Modified for LaTeX Templates by Vel (vel@LaTeXTemplates.com)
%
% License:
% The MIT License (see included LICENSE file)
%
%%%%%%%%%%%%%%%%%%%%%%%%%%%%%%%%%%%%%%%%%

%----------------------------------------------------------------------------------------
%	PACKAGES AND OTHER DOCUMENT CONFIGURATIONS
%----------------------------------------------------------------------------------------

\documentclass[9pt]{developercv} % Default font size, values from 8-12pt are recommended
\usepackage{graphicx}
\graphicspath{ {images/} }

%----------------------------------------------------------------------------------------

\begin{document}

%----------------------------------------------------------------------------------------
%	TITLE AND CONTACT INFORMATION
%----------------------------------------------------------------------------------------

\begin{minipage}[t]{0.45\textwidth} % 45% of the page width for name
	\vspace{-\baselineskip} % Required for vertically aligning minipages

	% If your name is very short, use just one of the lines below
	% If your name is very long, reduce the font size or make the minipage wider and reduce the others proportionately
	\colorbox{black}{{\HUGE\textcolor{white}{\textbf{\MakeUppercase{Nikita}}}}} % First name

	\colorbox{black}{{\HUGE\textcolor{white}{\textbf{\MakeUppercase{Razdobreev}}}}} % Last name

	\vspace{6pt}

	{\huge Gameplay Programmer} % Career or current job title

	\vspace{6pt}
	\includegraphics[scale=0.05]{photo}
	\vspace{6pt}
\end{minipage}
\begin{minipage}[t]{0.5\textwidth} % 27.5% of the page width for the first row of icons
	\vspace{-\baselineskip} % Required for vertically aligning minipages

	% The first parameter is the FontAwesome icon name, the second is the box size and the third is the text
	% Other icons can be found by referring to fontawesome.pdf (supplied with the template) and using the word after \fa in the command for the icon you want
	\icon{MapMarker}{12}{Saint-Petersburg, Russia}\\
	\icon{Phone}{12}{+7 (911) 844-73-54}\\
	\icon{At}{12}{\href{mailto:nikitarazdobreev94@gmail.com}{nikitarazdobreev94@gmail.com}}\\
	\icon{At}{12}{\href{mailto:nikita.razdobreev@protonmail.com}{nikita.razdobreev@protonmail.com}}\\
	\icon{Skype}{12}{nikita.razdobreev@protonmail.com}\\
	\icon{Linkedin}{12}{https://www.linkedin.com/in/razdobreev-nikita-0785b956}
	\icon{Github}{12}{\href{https://github.com/yahuaru}{github.com/yahuaru}}\\
	\icon{Twitter}{12}{\href{https://twitter.com/yahuaru}{yahuaru}}\\

\end{minipage}

\vspace{0.5cm}

%----------------------------------------------------------------------------------------
%	INTRODUCTION, SKILLS AND TECHNOLOGIES
%----------------------------------------------------------------------------------------

\cvsect{Skills}

\begin{center}
	\begin{tabular}{p{0.2\textwidth}p{0.2\textwidth}}
		\textbf{Languges:}     &              \\
		\textbf{Python}        & Expert       \\
		\textbf{C++}           & Intermediate \\
		\textbf{C\#}           & Intermediate \\
		\textbf{SQL}           & Intermediate \\
		                       &              \\
		\textbf{Engines:}      &              \\
		\textbf{Unreal Engine} & Intermediate \\
		\textbf{Unity}         & Beginner     \\
		                       &              \\
		\textbf{VCS:}          &              \\
		\textbf{SVN}           & Expert       \\
		\textbf{Git}           & Expert       \\
	\end{tabular}
\end{center}

%----------------------------------------------------------------------------------------
%	EXPERIENCE
%----------------------------------------------------------------------------------------

\cvsect{Experience}

\begin{entrylist}
	\entry
	{Apr 2019 -- Jan 2022}
	{Gameplay Programmer}
	{\href{https://na.wargaming.net/en/games/wowsconsole}{Wargaming Saint-Petersburg}}
	{

		Was responsible for support and development of "World of Warships: Legends"
		\begin{itemize}
			\item Implemented new core and meta gameplay features \\for online game with large user base
			\item Supported, fixed and refactored legacy code
			\item Developed new features for Server Core with decentralized architecture
			\item Integrated Server logic with RESTful Web Services
			\item Optimized memory, network and cpu workload for better performance on 8th generation of consoles
			\item Worked closely with 8th and 9th generation of consoles
			\item Developed editor toolset for gamedesigners
			\item Collaboration with a global team
		\end{itemize}
		\texttt{Python}\slashsep\texttt{Cython}\slashsep\texttt{C++}\slashsep\texttt{SQL}\slashsep\texttt{Docker}\slashsep\texttt{GIT}\slashsep\texttt{SVN}}
	\entry
	{Jan 2018 -- Mar 2019}
	{Gameplay Programmer}
	{\href{https://www.thepsycho.net/}{The Psycho®}}
	{

	Was responsible for gameplay programming and deployment of VR games
	\begin{itemize}
		\item Developed local multiplayer games for full body VR devices
		\item Implemented core gameplay features such as multiplayer shooting and AI for enemies
		\item Optimized cpu workload and content for best performance on VR devices
		\item Created workflow pipelines for content creators and VR tracking developers
		\item Levels and scenarios scripting
	\end{itemize}
	{\texttt{C++}\slashsep\texttt{Unreal Engine 4}\slashsep\texttt{Blueprint}\slashsep\texttt{GIT}\slashsep\texttt{SVN}}
	}
	\entry
	{Oct 2016 -- Nov 2017}
	{Junior Gameplay Programmer}
	{OctoBox interactive}
	{

	Was responsible for gameplay prototyping for unreleased Unity card game:
	\begin{itemize}
		\item Core gameplay and network programming for unreleased multiplayer card game on Unity
		\item Automation testing and TDD
	\end{itemize}
	{\texttt{C\#}\slashsep\texttt{Unity}\slashsep\texttt{Python}\slashsep\texttt{GIT}}
	}
	\entry
	{Feb 2016 -- Jul 2016}
	{Junior Gameplay Programmer}
	{\href{http://lazybeargames.com/}{Lazy Bear Games}}
	{

	Support for released Unity games
	\begin{itemize}
		\item QA and bug fixing for legacy projects on Unity
		\item Data analysis of player behavior for gameplay balancing
	\end{itemize}
	{\texttt{C\#}\slashsep\texttt{Unity}\slashsep\texttt{Python}}
	}
	\entry
	{Jan 2015 -- Nov 2015}
	{Junior Gameplay Programmer}
	{Islania}
	{

	Gameplay and UI prototyping for unreleased MMO \\ {\texttt{C++}\slashsep\texttt{Unreal Engine 4}\slashsep\texttt{Blueprint}}
	}
\end{entrylist}

%----------------------------------------------------------------------------------------
%	EDUCATION
%----------------------------------------------------------------------------------------

\cvsect{Education}

\begin{entrylist}
	\entry
	{2012-2016}
	{Bachelor's degree in Applied Mathematics \& Computer Science \\ \footnotesize{Saint Petersburg Electrotechnical University "LETI"} }
	{}
	{}
\end{entrylist}

%----------------------------------------------------------------------------------------
%	ADDITIONAL INFORMATION
%----------------------------------------------------------------------------------------

\begin{minipage}[t]{0.5\textwidth}
	\vspace{-\baselineskip} % Required for vertically aligning minipages

	\cvsect{Languages}

	\textbf{English} - Upper intermediate (B2)\\
	\textbf{Russian} - native
\end{minipage}
\hfill

%----------------------------------------------------------------------------------------

\end{document}
